\renewcommand{\Pr}{\operatorname{Pr}}

%%%%%%%%%%%%%%%%%%%%%%%%%%%%%%%%%%%%%%%%%%%%%%%%%%%%%%%%%%%%%%%%%%%%%%%%%%%%%%%
% Lossy compression and rate--distortion
%%%%%%%%%%%%%%%%%%%%%%%%%%%%%%%%%%%%%%%%%%%%%%%%%%%%%%%%%%%%%%%%%%%%%%%%%%%%%%%

\subsection{Setup: Rate--Distortion Tradeoff}
\label{subsec:lossy-setup}

For real-valued sources, perfect reconstruction with a finite number of bits is impossible.
A \emph{lossy} code trades rate (bits or nats per source symbol) for fidelity measured by a distortion function.
The central question is:
\begin{center}
\emph{How many bits per sample are needed to describe a random object up to distortion $D$?}
\end{center}

\paragraph{Single-letter distortion.}
Let $X\sim P_X$ take values in $\cX$ (often $\cX\subseteq\fR^d$).
Choose a distortion function $d:\cX\times\widehat{\cX}\to\fR_{\ge 0}$.
For mean-squared error (MSE),
\[
    d(x,\hat x)=(x-\hat x)^2,
\]
and the distortion level is $D=\E{(X-\widehat X)^2}$.
For binary sources, a common choice is Hamming distortion
\[
    d_H(x,\hat x)=\I{x\neq \hat x}.
\]
For model compression and sparse representations, an important distortion is absolute error
\[
    d_1(x,\hat x)=|x-\hat x|.
\]

\paragraph{Block codes.}
An $(n,R)$ lossy code consists of an encoder and decoder
\[
    f:\cX^n\to \{1,\dots,2^{nR}\},
    \qquad
    g:\{1,\dots,2^{nR}\}\to \widehat{\cX}^n,
\]
where $R$ is measured in bits per source symbol.
Equivalently, if $R$ is measured in nats, the message set has size $e^{nR}$.
The reconstruction is
\[
    \widehat X^n=g(f(X^n)).
\]
We measure average per-letter distortion by
\begin{equation}
    d^{(n)}(x^n,\hat x^n)
    \coloneqq \frac{1}{n}\sum_{i=1}^n d(x_i,\hat x_i),
    \qquad
    \E{d^{(n)}(X^n,\widehat X^n)}.
\end{equation}

\begin{definition}[Achievable $(R,D)$]
A pair $(R,D)$ is \emph{achievable} if for every $\varepsilon>0$ and all sufficiently large $n$, there exists an $(n,R+\varepsilon)$ code such that
\[
    \E{d^{(n)}(X^n,\widehat X^n)}\le D+\varepsilon.
\]
\end{definition}

\begin{definition}[Rate--distortion function]
The operational rate--distortion function is
\begin{equation}
    R(D)\coloneqq \inf\{R:\ (R,D)\text{ achievable}\}.
\end{equation}
\end{definition}

\paragraph{Interpretation.}
A lossless code must distinguish essentially all typical source sequences.
A lossy code only needs to distinguish \emph{distortion balls}: two source sequences can share the same description if the same reconstruction is close enough to both.
Rate--distortion theory asks for the smallest number of distortion balls needed to cover the typical set.

\subsection{Scalar Quantization: A Concrete Example}
\label{subsec:scalar-quant}

As a warm-up, consider a scalar Gaussian source (e.g., heights).
Assume $X\sim\cN(\mu,\sigma^2)$ with $\mu=170$ and $\sigma^2=10$.
We want to store $X$ using $R$ bits per sample.

\paragraph{Rate $R=0$.}
If we output a constant $\widehat X=\mu$, then
\[
    D(0)=\Var{X}=\sigma^2.
\]
The best zero-bit description is the mean, because the mean minimizes squared error among all constants.

\paragraph{Rate $R=1$ (one threshold).}
A 1-bit quantizer partitions at $\mu$:
\[
    f(x)=\I{x>\mu}.
\]
Given the index, the optimal reconstruction is the conditional mean,
\[
    \hat x_0=\E{X\mid X\le \mu},\qquad
    \hat x_1=\E{X\mid X>\mu}.
\]
For a Gaussian, this yields
\begin{equation}
    D(1)=\sigma^2\Bigl(1-\frac{2}{\pi}\Bigr)\approx 0.3634\,\sigma^2 \approx 3.63.
\end{equation}
The threshold does not store the precise value of $X$; it only tells us which half-line $X$ fell into.
The decoder then uses the best representative for that cell.

\paragraph{Rate $R=2$ (three thresholds).}
With four cells (e.g., thresholds $165,170,175$) the distortion improves, but the choice of thresholds is not obviously optimal.
The optimal thresholds and representatives depend on the source distribution and the distortion measure.

\subsection{Lloyd--Max Algorithm}
\label{subsec:lloyd-max}

For a fixed number of reconstruction levels, the Lloyd--Max algorithm alternates between:
\begin{enumerate}[leftmargin=2em]
    \item \textbf{Centroid update:} with thresholds $\{\tau_i\}$ fixed,
    \[
        \hat x_i \leftarrow \E{X\mid \tau_i\le X<\tau_{i+1}}.
    \]

    \item \textbf{Boundary update:} with representatives $\{\hat x_i\}$ fixed,
    \[
        \tau_{i+1}\leftarrow (\hat x_i+\hat x_{i+1})/2
    \]
    for squared-error distortion.
\end{enumerate}
Each step decreases distortion and the algorithm converges to a local optimum.
The important point is that quantization is not only about choosing labels; it is about jointly choosing the partition and the reproduction values.

\begin{figure}[ht]
\centering
    \begin{subfigure}{0.40\textwidth}
        \includegraphics[width=\linewidth]{scribe19_image2a.png}
        \caption{Fix thresholds, update representatives}
        \label{fig:lm-a}
    \end{subfigure}
    \hfill
    \begin{subfigure}[b]{0.08\textwidth}
        \includegraphics[width=\linewidth]{scribe19_arrow_both_directions.png}
    \end{subfigure}
    \hfill
    \begin{subfigure}[b]{0.40\textwidth}
        \includegraphics[width=\linewidth]{scribe19_image2b.png}
        \caption{Fix representatives, update thresholds}
        \label{fig:lm-b}
    \end{subfigure}
\caption{Alternating minimization in Lloyd--Max quantization.}
\label{fig:lm}
\end{figure}

\subsection{Vector Quantization}
\label{subsec:vq}

Scalar quantization treats each coordinate independently.
Vector quantization (VQ) quantizes blocks $X^n$ jointly:
\[
  Q^{(n)} : \fR^{n} \longrightarrow \{1,\dots,2^{nR}\}.
\]
Even for independent components, VQ can be strictly better because it can shape quantization regions to match the source geometry.
In high dimension, most probability mass is concentrated in a thin typical shell, and good vector quantizers place reconstruction points to cover this shell efficiently.

\begin{figure}[ht]
  \centering
  \includegraphics[width=.34\textwidth]{scribe19_image3.png}
  \caption{Example 2-D source.}
  \label{fig:vq-source}
\end{figure}

\begin{figure}[ht]
  \centering
  \begin{subfigure}{0.41\textwidth}
    \includegraphics[width=\linewidth]{scribe19_image4a.png}
    \caption{Scalar quantization (rectangles).}
    \label{fig:vq-scalar}
  \end{subfigure}\hspace{8mm}
  \begin{subfigure}{0.38\textwidth}
    \includegraphics[width=\linewidth]{scribe19_image4b.png}
    \caption{Vector quantization (adapted regions).}
    \label{fig:vq-vector}
  \end{subfigure}
  \caption{Scalar vs. vector quantization for a 2-D uniform source.}
  \label{fig:vq}
\end{figure}

\begin{figure}[ht]
  \centering
  \includegraphics[width=0.30\textwidth]{scribe19_Figure3.jpg}
  \caption{Rectangular cells incur extra MSE compared to more circular cells.}
  \label{fig:vq-distortion}
\end{figure}

\subsection{Information-Theoretic Rate--Distortion Function}
\label{sec:rd-info}

The operational definition of $R(D)$ admits an exact single-letter characterization.

\begin{definition}[Information rate--distortion function]
\label{def:rd-info}
For a memoryless source $X\sim P_X$ and single-letter distortion $d$, define
\begin{equation}
    R^{(I)}(D)
    \coloneqq
    \inf_{P_{\widehat X\mid X}:\,\E{d(X,\widehat X)}\le D}
    I(X;\widehat X).
\end{equation}
The infimum is over all conditional distributions (``test channels'') from $X$ to $\widehat X$.
\end{definition}

\begin{theorem}[Shannon's rate--distortion theorem]
\label{thm:shannon-rd}
For every memoryless source and single-letter distortion measure,
\begin{equation}
    R(D)=R^{(I)}(D).
\end{equation}
\end{theorem}

\paragraph{Why mutual information appears.}
A rate-$R$ code selects one out of $2^{nR}$ reconstructions.
The random index cannot carry more than $nR$ bits about the source block.
Thus the reconstruction cannot contain more than $nR$ bits of information about the source.
The theorem says that the exact per-symbol information needed to reproduce $X$ up to distortion $D$ is
\[
    I(X;\widehat X),
\]
minimized over all test channels whose average distortion is at most $D$.

\subsection{Random Coding: Why a Random Codebook Covers the Source}
\label{subsec:random-coding-intuition}

The achievability part of Theorem~\ref{thm:shannon-rd} is often the least intuitive part.
The core idea is surprisingly simple:
\begin{center}
\emph{Generate many independent fake reconstructions.  One of them will look as if it was generated jointly with the true source.}
\end{center}

Fix a test channel $P_{\widehat X\mid X}$ with joint distribution
\[
    P_{X,\widehat X}(x,\hat x)=P_X(x)P_{\widehat X\mid X}(\hat x\mid x)
\]
and distortion $\E{d(X,\widehat X)}\le D$.
Let $P_{\widehat X}$ be the marginal reproduction distribution.
For blocklength $n$, generate a random codebook
\begin{equation}
    \widehat X^n(1),\ldots,\widehat X^n(M),
    \qquad M=2^{nR},
\end{equation}
where each codeword is drawn i.i.d. according to $P_{\widehat X}$.
Given a source sequence $X^n$, the encoder searches for an index $m$ such that
\[
    (X^n,\widehat X^n(m))
\]
is jointly typical with respect to $P_{X,\widehat X}$.
If such an $m$ exists, the encoder sends $m$ and the decoder outputs $\widehat X^n(m)$.

\paragraph{The rare-event calculation.}
For a fixed source sequence $X^n$ that is typical under $P_X$, and for one independently generated codeword $\widehat X^n\sim P_{\widehat X}^{\otimes n}$,
\begin{equation}
    \Pr\bigl[(X^n,\widehat X^n)\text{ jointly typical}\bigr]
    \approx 2^{-nI(X;\widehat X)}.
\end{equation}
This is because joint typicality asks for the pair to behave as if it came from $P_{X,\widehat X}$ rather than from the independent product $P_XP_{\widehat X}$.
The ``penalty'' for forcing this dependence is exactly mutual information.
With $M=2^{nR}$ independent trials, the expected number of jointly typical codewords is approximately
\begin{equation}
    2^{nR}\cdot 2^{-nI(X;\widehat X)}=2^{n(R-I(X;\widehat X))}.
\end{equation}
If $R>I(X;\widehat X)$, this number grows exponentially.
Moreover,
\begin{equation}
    \Pr[\text{no jointly typical codeword}]
    \approx
    \left(1-2^{-nI(X;\widehat X)}\right)^{2^{nR}}
    \le \exp\{-2^{n(R-I(X;\widehat X))}\},
\end{equation}
which vanishes rapidly.
Thus a random codebook of rate slightly larger than $I(X;\widehat X)$ contains, with high probability, at least one reconstruction sequence that is jointly typical with $X^n$.

\paragraph{What ``as if jointly sampled'' means.}
The codewords were generated independently of $X^n$.
But the encoder does not choose a uniformly random codeword; it chooses a codeword satisfying a rare joint-typicality event.
Conditioned on this event, the selected pair $(X^n,\widehat X^n(m))$ has empirical frequencies close to $P_{X,\widehat X}$.
Therefore, it has distortion close to
\[
    \E{d(X,\widehat X)}.
\]
The selection step creates the desired statistical dependence.
This is the same phenomenon as drawing many independent strings and keeping the one whose empirical correlation with $X^n$ has the desired value.

\subsection{Visual Binary Example: Bernoulli Source and Hamming Distortion}
\label{subsec:binary-visual}

Let
\[
    X_i\overset{\text{i.i.d.}}{\sim}\mathrm{Bern}(1/2),
    \qquad
    d_H(x,\hat x)=\I{x\neq \hat x}.
\]
For any $0<D<1/2$, the optimal test channel can be written as
\begin{equation}
    X=\widehat X\oplus Z,
    \qquad
    \widehat X\sim \mathrm{Bern}(1/2),
    \qquad
    Z\sim \mathrm{Bern}(D),
\end{equation}
with $\widehat X$ independent of $Z$.
Equivalently, $\widehat X=X\oplus Z$.
The reproduction marginal is again Bernoulli$(1/2)$, but $X$ and $\widehat X$ are not independent: they differ in a fraction $D$ of coordinates.
Hence
\begin{equation}
    I(X;\widehat X)=H(X)-H(X\mid \widehat X)=1-\Hb{D}
    \qquad\text{bits/symbol},
\end{equation}
and the rate--distortion function is
\begin{equation}
R(D)=
\begin{cases}
1-\Hb{D}, & 0\le D\le 1/2,\\
0, & D\ge 1/2.
\end{cases}
\end{equation}

\paragraph{Independent codewords look unrelated.}
Suppose $n=24$ and $D=1/4$.
Then a good reconstruction should be within $nD=6$ bit flips of $X^n$.
Let
\[
X^n=\texttt{101000100001100010000100}.
\]
If we draw one random Bernoulli$(1/2)$ codeword, it typically differs from $X^n$ in about $12$ coordinates.
A few independently sampled candidates might look like this:
\begin{center}
\begin{tabular}{c|c|c}
index $m$ & candidate $\widehat X^n(m)$ & Hamming distance to $X^n$ \\
\hline
1 & \texttt{110111111001001010011011} & 14 \\
2 & \texttt{101110001011010000010011} & 12 \\
5 & \texttt{100000000001010000001000} & 7 \\
6 & \texttt{101000111001001011100000} & 7
\end{tabular}
\end{center}
None of these four candidates is close enough for distortion $D=1/4$.
This is expected: each candidate has the correct marginal distribution, but it is independent of $X^n$.

\paragraph{Many independent codewords contain a correlated-looking one.}
If we keep sampling more Bernoulli$(1/2)$ codewords, eventually one may land in the Hamming ball of radius $6$ around $X^n$.
For example, in a random codebook of size $64$, one possible candidate is
\begin{center}
\begin{tabular}{c|c|c}
index $m$ & candidate $\widehat X^n(m)$ & Hamming distance to $X^n$ \\
\hline
14 & \texttt{101010101000100111000000} & 6
\end{tabular}
\end{center}
The error pattern for this candidate is
\[
X^n\oplus \widehat X^n(14)
=	\texttt{000010001001000101000100},
\]
which has six ones, i.e., empirical crossover probability $6/24=1/4$.
Thus the selected pair behaves like
\[
    \widehat X_i = X_i\oplus Z_i,
    \qquad Z_i\sim\mathrm{Bern}(1/4),
\]
even though every candidate was generated independently before seeing $X^n$.
The encoder finds dependence by search.

\paragraph{Threshold rate.}
The volume of a Hamming ball of radius $nD$ is approximately
\[
    2^{n\Hb{D}}.
\]
Since the whole binary space has size $2^n$, a random Bernoulli codeword falls in the distortion-$D$ ball with probability approximately
\[
    2^{-n(1-\Hb{D})}.
\]
Therefore $2^{nR}$ random candidates will contain a close one with high probability when
\[
    R>1-\Hb{D}.
\]
For $D=1/4$,
\[
    1-\Hb{1/4}\approx 0.189\quad\text{bits/symbol}.
\]
Finite blocklengths need a safety margin, but asymptotically this is the exact threshold.

\subsection{Gaussian Source with Quadratic Distortion}
\label{subsec:gaussian-rd}

Let $X\sim\cN(0,\sigma^2)$ and $d(x,\hat x)=(x-\hat x)^2$.
Then
\begin{equation}
R(D)=
\begin{cases}
\tfrac12\log\Bigl(\dfrac{\sigma^2}{D}\Bigr), & 0<D\le \sigma^2,\\[1mm]
0, & D\ge \sigma^2,
\end{cases}
\label{eq:gaussian-rd}
\end{equation}
where $\log$ is natural logarithm, so the rate is in nats.
To express the rate in bits, divide the right-hand side by $\log 2$.

\subsubsection*{Achievability: backward Gaussian test channel}

Fix $0<D<\sigma^2$.
Let
\begin{equation}
    \widehat X\sim \cN(0,\sigma^2-D),
    \qquad
    Z\sim \cN(0,D),
    \qquad
    X=\widehat X+Z,
\end{equation}
with $\widehat X$ independent of $Z$.
Then $X\sim\cN(0,\sigma^2)$ and
\[
    \E{(X-\widehat X)^2}=\E{Z^2}=D.
\]
The mutual information is
\begin{align}
    I(X;\widehat X)
    &= h(X)-h(X\mid \widehat X) \\
    &= h(X)-h(Z) \\
    &= \frac12\log(2\pi e\sigma^2)-\frac12\log(2\pi eD) \\
    &= \frac12\log\frac{\sigma^2}{D}.
\end{align}
By the random-coding argument in Section~\ref{subsec:random-coding-intuition}, generate
\[
    e^{nR}
\]
codewords independently from $\cN(0,\sigma^2-D)^{\otimes n}$.
If $R>I(X;\widehat X)$, then with high probability there exists a codeword $\widehat X^n(m)$ such that
\[
    X^n\approx \widehat X^n(m)+Z^n,
    \qquad Z_i\overset{\text{i.i.d.}}{\sim}\cN(0,D),
\]
in the sense of joint typicality.
The resulting empirical squared error is close to $D$.
This proves achievability of any rate larger than the expression in~\eqref{eq:gaussian-rd}.

\paragraph{Forward form of the same test channel.}
The above test channel is called ``backward'' because it writes $X$ as a noisy version of $\widehat X$.
Equivalently, the conditional distribution of $\widehat X$ given $X=x$ is Gaussian:
\begin{equation}
    \widehat X\mid X=x
    \sim
    \cN\left(\frac{\sigma^2-D}{\sigma^2}x,
    \frac{D(\sigma^2-D)}{\sigma^2}\right).
\end{equation}
This is the best linear shrinkage of $X$ plus additional randomness.
The reconstruction has smaller variance than the source, because a lossy description cannot reproduce all the source energy.

\subsubsection*{Converse: no code can beat the Gaussian formula}

Let $M=f(X^n)$ be the message of any $(n,R)$ code and let $\widehat X^n=g(M)$.
Assume
\[
    \frac1n\E{\|X^n-\widehat X^n\|_2^2}\le D.
\]
Since $M$ takes at most $e^{nR}$ values,
\begin{equation}
    nR\ge H(M)\ge I(X^n;M)\ge I(X^n;\widehat X^n).
\end{equation}
Now let $E^n=X^n-\widehat X^n$.
Then
\begin{align}
I(X^n;\widehat X^n)
&= h(X^n)-h(X^n\mid \widehat X^n) \\
&= h(X^n)-h(X^n-\widehat X^n\mid \widehat X^n) \\
&\ge h(X^n)-h(E^n).
\end{align}
The Gaussian source has
\[
    h(X^n)=\frac n2\log(2\pi e\sigma^2).
\]
Among all random vectors with second moment constraint
\[
    \E{\|E^n\|_2^2}\le nD,
\]
the entropy is maximized by an isotropic Gaussian with covariance $DI$.
Therefore
\begin{equation}
    h(E^n)\le \frac n2\log(2\pi eD).
\end{equation}
Combining the inequalities gives
\begin{equation}
    nR\ge \frac n2\log\frac{\sigma^2}{D},
\end{equation}
so
\[
    R\ge \frac12\log\frac{\sigma^2}{D}.
\]
This proves the converse.

\paragraph{Geometric picture.}
In $\fR^n$, a typical Gaussian source block lies near a sphere of radius $\sqrt{n\sigma^2}$.
A distortion constraint $D$ means each codeword covers a ball of radius $\sqrt{nD}$.
A crude volume comparison gives
\[
    \frac{\text{volume of source shell}}{\text{volume of one distortion ball}}
    \approx
    \left(\frac{\sqrt{n\sigma^2}}{\sqrt{nD}}\right)^n
    =\left(\frac{\sigma^2}{D}\right)^{n/2}.
\]
Taking logarithms and dividing by $n$ gives
\[
    R\approx \frac12\log\frac{\sigma^2}{D}.
\]
The entropy proof above is the rigorous version of this geometric argument.

\begin{figure}[h]
\centering
\includegraphics[width=0.58\textwidth]{scribe27_sphere.jpg}
\caption{High-dimensional picture for Gaussian rate--distortion: cover a large sphere (source) by many small balls (distortion).}
\label{fig:gauss-geometry}
\end{figure}

\subsection{Shannon Lower Bound and Why Gaussian Sources Are Worst Case}
\label{subsec:slb-gaussian-worst}

The Gaussian formula is not only an example.
Under squared-error distortion, Gaussian sources are the hardest to compress among sources with the same variance.
The main tool is the Shannon lower bound.

\begin{theorem}[Shannon lower bound under squared error]
\label{thm:slb}
Let $X$ be a real-valued source with finite differential entropy $h(X)$.
For squared-error distortion,
\begin{equation}
    R_X(D)\ge h(X)-\frac12\log(2\pi eD).
    \label{eq:slb-scalar}
\end{equation}
More generally, for $X\in\fR^d$ and per-coordinate distortion
$\frac1d\E{\|X-\widehat X\|_2^2}\le D$,
\begin{equation}
    R_X(D)\ge h(X)-\frac d2\log(2\pi eD).
    \label{eq:slb-vector}
\end{equation}
\end{theorem}

\begin{proof}
For any reconstruction $\widehat X$ satisfying $\E{(X-\widehat X)^2}\le D$,
let $E=X-\widehat X$.
Then
\begin{align}
I(X;\widehat X)
&=h(X)-h(X\mid \widehat X) \\
&=h(X)-h(X-\widehat X\mid \widehat X) \\
&\ge h(X)-h(X-\widehat X) \\
&=h(X)-h(E).
\end{align}
Among all scalar random variables with variance at most $D$, the Gaussian has the largest entropy.
Hence
\[
    h(E)\le \frac12\log(2\pi eD),
\]
which proves~\eqref{eq:slb-scalar}.
For the vector case, use
\[
    h(E^d)\le \frac12\log\bigl((2\pi e)^d\det K_E\bigr),
\]
where $K_E$ is the covariance matrix of $E^d$.
If $\mathrm{tr}(K_E)\le dD$, then by the arithmetic--geometric mean inequality,
\[
    \det K_E\le \left(\frac{\mathrm{tr}(K_E)}{d}\right)^d\le D^d.
\]
Thus $h(E^d)\le \frac d2\log(2\pi eD)$.
\end{proof}

\paragraph{Gaussian achieves the Shannon lower bound.}
If $X\sim\cN(0,\sigma^2)$, then
\[
    h(X)=\frac12\log(2\pi e\sigma^2).
\]
Substituting this into the Shannon lower bound gives
\[
    R_X(D)\ge \frac12\log\frac{\sigma^2}{D}.
\]
The backward Gaussian test channel achieves equality.
Thus the Shannon lower bound is tight for Gaussian sources.

\paragraph{Worst-case interpretation.}
Suppose a source has variance $\sigma^2$ and finite second moment.
It is always possible to cover the ball of radius approximately $\sqrt{n\sigma^2}$ by distortion balls of radius approximately $\sqrt{nD}$, yielding a universal lossy code with rate close to
\[
    \frac12\log\frac{\sigma^2}{D}.
\]
Therefore
\begin{equation}
    R_X(D)\le \frac12\log\frac{\sigma^2}{D}
\end{equation}
for any source with variance $\sigma^2$ (up to the usual asymptotic $\varepsilon$ terms).
The Gaussian source attains this upper bound exactly.
In this precise sense, among all sources with the same variance, the Gaussian is the least compressible under MSE distortion.

\paragraph{Intuition.}
Gaussian distributions have the largest differential entropy among all distributions with a fixed variance.
They spread probability mass as diffusely as possible.
A diffuse source has a large typical set, so more codewords are needed to cover it at the same distortion level.
Non-Gaussian sources with the same variance have more structure and can usually be compressed at a lower rate.

\subsection{Laplacian Source with Absolute Distortion}
\label{subsec:laplace-l1-rd}

Let
\[
    X\sim\mathrm{Lap}(\lambda),
    \qquad
    f_X(x)=\frac{\lambda}{2}e^{-\lambda |x|},
\]
so $\E{|X|}=1/\lambda$.
Use absolute distortion
\[
    d_1(x,\hat x)=|x-\hat x|.
\]
Then the rate--distortion function is
\begin{equation}
R(D)=
\begin{cases}
-
\log(\lambda D), & 0<D\le 1/\lambda,\\[1mm]
0, & D\ge 1/\lambda.
\end{cases}
\label{eq:laplace-l1-rd}
\end{equation}
This example is central for model pruning, because normalized neural-network weights are often empirically close to Laplacian, and the $\ell_1$ distortion controls output perturbation in a useful way.

\subsubsection*{Converse}

Let $\widehat X$ be any reconstruction satisfying
\[
    \E{|X-\widehat X|}\le D.
\]
Consider the auxiliary conditional density
\begin{equation}
    q_{X\mid \widehat X}(x\mid \hat x)
    =\frac{1}{2D}\exp\left(-\frac{|x-\hat x|}{D}\right),
\end{equation}
i.e., a Laplace density centered at $\hat x$ with scale $D$.
Using nonnegativity of relative entropy,
\begin{align}
I(X;\widehat X)
&=D_{\mathrm{KL}}(P_{X\mid \widehat X}\|P_X\mid P_{\widehat X}) \\
&\ge
\E{\log\frac{q_{X\mid \widehat X}(X\mid \widehat X)}{p_X(X)}}.
\end{align}
Now compute the expectation:
\begin{align}
\E{\log\frac{q_{X\mid \widehat X}(X\mid \widehat X)}{p_X(X)}}
&=
\E{\log\frac{(2D)^{-1}e^{-|X-\widehat X|/D}}{(\lambda/2)e^{-\lambda |X|}}} \\
&= -\log(\lambda D)-\frac{1}{D}\E{|X-\widehat X|}+\lambda\E{|X|} \\
&\ge -\log(\lambda D)-1+1 \\
&= -\log(\lambda D).
\end{align}
Thus $I(X;\widehat X)\ge -\log(\lambda D)$ for every admissible reconstruction.

\subsubsection*{Achievability}

We now construct a test channel that achieves the lower bound.
Let $N\sim\mathrm{Lap}(1/D)$ have density
\[
    f_N(n)=\frac{1}{2D}e^{-|n|/D}.
\]
Let $V$ be independent of $N$ with distribution
\begin{equation}
    P_V
    =\lambda^2D^2\,\delta_0
    +(1-\lambda^2D^2)\,\mathrm{Lap}(\lambda),
    \label{eq:laplace-reproduction-mixture}
\end{equation}
where $\delta_0$ is a point mass at zero.
For $0<D\le 1/\lambda$, the mixture weight $\lambda^2D^2$ is in $[0,1]$.
One can verify by characteristic functions that
\begin{equation}
    X\coloneqq V+N \sim \mathrm{Lap}(\lambda).
\end{equation}
Indeed,
\[
\phi_N(t)=\frac{1}{1+D^2t^2},
\qquad
\phi_X(t)=\frac{\lambda^2}{\lambda^2+t^2},
\]
and the mixture distribution in~\eqref{eq:laplace-reproduction-mixture} has characteristic function
\[
\phi_V(t)=\lambda^2D^2+(1-\lambda^2D^2)\frac{\lambda^2}{\lambda^2+t^2},
\]
so $\phi_V(t)\phi_N(t)=\phi_X(t)$.
Taking $\widehat X=V$, we have
\[
    \E{|X-\widehat X|}=\E{|N|}=D.
\]
The conditional density of $X$ given $V=v$ is exactly the auxiliary density used in the converse, so equality holds and
\[
    I(X;V)=-\log(\lambda D).
\]

\paragraph{Why sparsity appears.}
The optimal reproduction distribution $P_V$ has an atom at zero with probability $\lambda^2D^2$.
Thus, at the rate--distortion limit, a positive fraction of reconstructed symbols are exactly zero.
This is not an artifact of a particular algorithm: it is a property of the optimal test channel.
For Laplacian sources under $\ell_1$ distortion, sparsity is information-theoretically natural.

\paragraph{Exponential one-sided version.}
For magnitudes $U=|X|$, the natural source is exponential:
\[
    f_U(u)=\lambda e^{-\lambda u},\qquad u\ge 0.
\]
With one-sided distortion
\[
    d(u,v)=
    \begin{cases}
    u-v, & u\ge v,\\
    \infty, & u<v,
    \end{cases}
\]
the rate--distortion function is again
\begin{equation}
    R(D)=-\log(\lambda D),\qquad 0<D\le 1/\lambda,
\end{equation}
and the optimal reconstruction has distribution
\begin{equation}
    P_V=\lambda D\,\delta_0+(1-\lambda D)\,\mathrm{Exp}(\lambda).
\end{equation}
This is useful for neural networks because it lets us protect sign bits separately and compress only magnitudes.

\subsection{Extreme-Value Theory as a Tool for Sparse Refinement}
\label{sec:evt}

Some practical rate--distortion schemes repeatedly identify unusually large coordinates of a residual vector.
Extreme-value theory tells us where these maxima typically lie.

\begin{example}[Max of exponentials]
If $Y_1,\dots,Y_n\overset{\text{i.i.d.}}{\sim}\mathrm{Exp}(\lambda)$ and $M_n=\max_i Y_i$, then
\begin{equation}
    M_n = \frac{\log n}{\lambda} + O_p(1).
\end{equation}
The fluctuations converge to a Gumbel distribution after centering by $\frac{1}{\lambda}\log n$.
\end{example}

\paragraph{Rule of thumb.}
For i.i.d. $X_i\sim\cN(0,\sigma^2)$,
\begin{equation}
    \max_i X_i = \sigma\sqrt{2\log n}+O_p(1).
\end{equation}
Thus, the value of the maximum is almost predictable from $n$ alone.
This is why an encoder can sometimes describe only the \emph{index} of an extreme coordinate and let the decoder infer a good reconstruction value.

\subsection{CROM: Rateless Lossy Compression via Extremes}
\label{subsec:crom}

Coding with Random Orthogonal Matrices (CROM) is a constructive lossy compression scheme built from two ideas:
\begin{enumerate}[leftmargin=2em]
    \item at very low rate, describe only the indices of a few extreme coordinates;
    \item after each sparse description, rotate the residual so that its coordinates again look approximately Gaussian.
\end{enumerate}
The result is an iterative, rateless compression method that achieves the Gaussian distortion--rate curve while using sparse messages at each step.

\subsubsection*{Zero-rate compression by the maximum}

Let $X^n\sim\cN(0,I_n)$.
The maximum coordinate satisfies
\[
    \max_{1\le i\le n}X_i\approx \sqrt{2\log n}.
\]
A near-zero-rate encoder can send only
\[
    m=\arg\max_i X_i.
\]
This costs $\log n$ nats total, or
\[
    R_n=\frac{\log n}{n}\to 0
\]
nats per symbol.
The decoder reconstructs a vector with a positive spike at the selected coordinate and a small negative offset elsewhere, for example
\begin{equation}
\widehat X_i=
\begin{cases}
\alpha_n, & i=m,\\
-\dfrac{\alpha_n}{n-1}, & i\neq m,
\end{cases}
\qquad \alpha_n\approx \sqrt{2\log n}.
\end{equation}
The negative offset keeps the reconstruction roughly mean-zero.
The distortion reduction is approximately
\[
    \frac{2\log n}{n}=2R_n.
\]
For a unit-variance Gaussian source,
\[
    D_G(R)=e^{-2R}=1-2R+o(R)
    \qquad (R\downarrow 0),
\]
so the zero-rate slope is exactly $-2$.
Thus this simple maximum-index scheme extracts the correct first-order distortion reduction per nat at zero rate.

\paragraph{Using several extremes.}
Instead of sending only the largest coordinate, the encoder may send the indices of the $k_n$ largest coordinates, often with $k_n=(\log n)^\beta$.
The rate is still vanishing:
\[
    R_n=\frac1n\log {n\choose k_n}\to 0,
\]
but the scheme can reduce distortion more robustly.
The decoder places equal positive mass on the selected coordinates and a compensating negative mass on the unselected coordinates.

\subsubsection*{Why rotation is needed}

After the first sparse reconstruction, the residual
\[
    X^n-\widehat X^n
\]
is no longer i.i.d. Gaussian.
Its largest coordinate has been partially removed, and the coordinates have a complicated dependence on the selected index.
If we immediately apply the same maximum-index idea again, the residual may not have the clean extreme-value behavior we need.

CROM solves this by multiplying the residual by a random orthogonal matrix $A$.
For any fixed residual vector $r^n$, the rotated vector $Ar^n$ is uniformly distributed on the sphere of radius $\|r^n\|_2$.
In high dimension, a uniformly random point on a sphere has coordinates that look nearly i.i.d. Gaussian with variance $\|r^n\|_2^2/n$.
Thus rotation ``Gaussianizes'' the residual enough to reuse the extreme-value compressor.

\subsubsection*{Iterative CROM idea}

Set $X^{(1)}=A_1X^n$.
At iteration $t$:
\begin{enumerate}[leftmargin=2em]
    \item Find the set $m^{(t)}$ of the $k_n$ largest coordinates of $X^{(t)}$.
    \item Encode only this index set.
    \item Let $U^{(t)}$ be a sparse unit vector supported on $m^{(t)}$, with a compensating offset on the remaining coordinates.
    \item Subtract a scaled version $\alpha_t U^{(t)}$.
    \item Rotate the new residual:
    \[
        X^{(t+1)}=A_{t+1}\bigl(X^{(t)}-\alpha_tU^{(t)}\bigr).
    \]
\end{enumerate}
The decoder receives the index sets $m^{(1)},m^{(2)},\ldots$ and reconstructs by undoing the rotations and summing the sparse refinements.
After $L$ iterations, the reconstruction has the form
\begin{equation}
    \widehat X^{(L)}
    =\alpha_1A_1^TU^{(1)}+
    \alpha_2A_1^TA_2^TU^{(2)}+
    \cdots+
    \alpha_LA_1^T\cdots A_L^TU^{(L)}.
\end{equation}

\paragraph{Ratelessness.}
The decoder can stop after receiving the first $i$ refinement messages.
At that point, it has a reconstruction at rate approximately $iR/L$.
CROM is designed so that every prefix of the message stream lies near the Gaussian distortion--rate curve
\[
    D_G(R)=\sigma^2e^{-2R}.
\]
Thus the encoder need not know the final desired rate in advance: the bitstream is progressively useful.
This is analogous to rateless channel coding, except that the decoder needs the initial prefix of the stream, not an arbitrary subset of packets.

\paragraph{Takeaway.}
Random coding proves existence by drawing an exponentially large codebook.
CROM gives a more structured construction: each refinement sends only sparse extreme-coordinate information after a random rotation.
The rotation step is what makes the simple zero-rate extreme-value idea reusable.

\subsection{Application: Compressing Neural Networks}
\label{sec:nn-compress}

Modern neural network (NN) compression often starts by choosing a distortion metric on weights that controls the induced change in the network output.
Rate--distortion theory then provides a principled limit for how well such weights can be compressed.
This section explains why pruning appears naturally from the Laplacian $\ell_1$ rate--distortion function.

\subsubsection*{A distortion metric that controls output perturbation}

Consider a $d$-layer fully-connected network
\[
    f(x;w)=W^{(d)}\sigma(W^{(d-1)}\sigma(\cdots \sigma(W^{(1)}x)\cdots))
\]
with a 1-Lipschitz nonlinearity $\sigma$ such that $\sigma(0)=0$ (e.g., ReLU).
Let $\widehat w$ be the compressed weights.
A useful bound is the following.

\begin{theorem}[Weight distortion implies output distortion]
\label{thm:nn-l1-bound}
Assume $\|\widehat W^{(\ell)}\|_1\le \|W^{(\ell)}\|_1$ for all layers.
Then
\begin{equation}
\sup_{\|x\|_1\le 1}\,\|f(x;w)-f(x;\widehat w)\|_1
\le
\Bigl(\prod_{k=1}^{d}\|W^{(k)}\|_1\Bigr)
\Bigl(\sum_{\ell=1}^{d}\frac{\|W^{(\ell)}-\widehat W^{(\ell)}\|_1}{\|W^{(\ell)}\|_1}\Bigr).
\label{eq:nn-output-bound}
\end{equation}
\end{theorem}

\paragraph{Proof idea.}
Write the network output difference as a telescoping perturbation over layers.
At the last layer,
\begin{align}
&\|W^{(d)}\sigma(h)-\widehat W^{(d)}\sigma(\widehat h)\|_1 \\
&\quad\le
\|(W^{(d)}-\widehat W^{(d)})\sigma(h)\|_1
+
\|\widehat W^{(d)}(\sigma(h)-\sigma(\widehat h))\|_1.
\end{align}
The first term is controlled by $\|W^{(d)}-\widehat W^{(d)}\|_1$.
The second term propagates the perturbation from earlier layers, multiplied by $\|\widehat W^{(d)}\|_1\le \|W^{(d)}\|_1$.
Repeating this argument layer by layer gives~\eqref{eq:nn-output-bound}.

The quantity
\begin{equation}
    \sum_{\ell=1}^{d}\frac{\|W^{(\ell)}-\widehat W^{(\ell)}\|_1}{\|W^{(\ell)}\|_1}
\end{equation}
is therefore a meaningful normalized weight distortion.
For a single vector of normalized weights $u=(u_1,\ldots,u_n)$, this motivates the distortion
\begin{equation}
    d(u,\widehat u)=\frac1n\sum_{i=1}^n |u_i-\widehat u_i|.
\end{equation}

\subsubsection*{Why Laplacian weights lead to pruning}

Empirically, many pretrained NN weights, after layerwise normalization, are well approximated by a Laplacian distribution.
If
\[
    U\sim \mathrm{Lap}(\lambda),
    \qquad
    d(u,\widehat u)=|u-\widehat u|,
\]
then Section~\ref{subsec:laplace-l1-rd} shows that the optimal reproduction distribution is
\[
    P_{\widehat U}
    =\lambda^2D^2\delta_0+(1-\lambda^2D^2)\mathrm{Lap}(\lambda).
\]
This contains an atom at zero.
Therefore, an information-theoretically optimal compressed model is not merely quantized; it is sparse.
In other words:
\begin{center}
\emph{For Laplacian-like weights under $\ell_1$ distortion, pruning is not just a heuristic.  It is part of the optimal reproduction law.}
\end{center}
A compression method that never outputs exact zeros cannot match this optimal marginal distribution and is generally suboptimal in the rate--distortion sense.

\subsubsection*{Sign bits and magnitude compression}

Changing the sign of a weight can have a large effect on the network output.
A practical approach is therefore:
\begin{enumerate}[leftmargin=2em]
    \item store or protect the sign pattern;
    \item compress the magnitudes $|w_i|$.
\end{enumerate}
If magnitudes are modeled as exponential, the one-sided exponential rate--distortion function again has a sparse optimal reconstruction:
\[
    P_{\widehat U}=\lambda D\delta_0+(1-\lambda D)\mathrm{Exp}(\lambda).
\]
This gives the same qualitative conclusion: many reconstructed magnitudes should be zero.

\subsubsection*{Successive Refinement for Pruning (SuRP)}

A random codebook that exactly achieves the rate--distortion function would have exponentially many codewords and is not practical for a neural network with millions of weights.
The SuRP idea is to exploit successive refinement and extreme values.
The reconstruction begins as the all-zero vector and gradually adds back large coordinates.

For the symmetric Laplacian version, initialize $U^{(1)}=u$ and estimate the Laplacian parameter $\lambda_1$.
At iteration $t$:
\begin{enumerate}[leftmargin=2em]
    \item Find an index $i$ with large positive residual and an index $j$ with large negative residual:
    \[
        U_i^{(t)}\ge \frac1{\lambda_t}\log\frac{n}{2\beta},
        \qquad
        U_j^{(t)}\le -\frac1{\lambda_t}\log\frac{n}{2\beta}.
    \]
    \item Send the pair of indices $(i,j)$.
    \item The decoder adds a positive spike at $i$ and a negative spike at $j$:
    \[
        \widehat U_i \leftarrow \widehat U_i+\frac1{\lambda_t}\log\frac{n}{2\beta},
        \qquad
        \widehat U_j \leftarrow \widehat U_j-\frac1{\lambda_t}\log\frac{n}{2\beta}.
    \]
    \item The encoder subtracts the same spikes from the residual.
    \item Update the effective Laplacian parameter for the next residual.
\end{enumerate}
For the magnitude-only exponential version, the encoder sends one index whose magnitude residual exceeds
\[
    \frac1{\lambda_t}\log\frac{n}{\beta},
\]
and the decoder adds that positive amount to the reconstructed magnitude.
The signs are then restored separately.

\paragraph{Relation to magnitude pruning.}
Extreme-value theory says the largest Laplacian or exponential samples are near a logarithmic threshold.
Thus SuRP repeatedly communicates large residual coordinates.
Magnitude pruning can be viewed as a non-iterative cousin of this idea: keep large weights and remove small weights.
SuRP differs in that the threshold and refinement rule come from the rate--distortion model rather than from an externally chosen pruning score.

\paragraph{Zero-rate optimality.}
Each SuRP iteration communicates only one or two indices, so the rate per weight is on the order of
\[
    \frac{\log n}{n}\to 0.
\]
For the Laplacian source with $\ell_1$ distortion, the distortion--rate function is
\[
    D(R)=\frac1\lambda e^{-R},
\]
whose zero-rate slope is
\[
    D'(0)=-\frac1\lambda.
\]
A single SuRP refinement reduces distortion at this optimal first-order slope, under the usual choice of threshold parameter (for example, $\beta$ growing like $\log n$).
Thus SuRP is a zero-rate optimal refinement step.
Repeated refinements build a sparse reconstruction progressively.

\subsubsection*{Information-theoretic justification of pruning}

The argument can be summarized in four steps.
\begin{enumerate}[leftmargin=2em]
    \item \textbf{Output perturbation can be bounded by normalized $\ell_1$ weight distortion.}
    This connects model compression to a classical lossy source coding problem.

    \item \textbf{Normalized weights are often close to Laplacian.}
    This suggests using the Laplacian source model rather than a Gaussian model for weights.

    \item \textbf{The Laplacian $\ell_1$ rate--distortion function has a sparse optimal reproduction distribution.}
    The atom at zero is the mathematical signature of pruning.

    \item \textbf{A practical refinement algorithm communicates extreme residual coordinates.}
    This produces a sparse model and has a zero-rate optimality guarantee.
\end{enumerate}
Therefore, pruning is not only an empirical trick for overparameterized networks.
Under a Laplacian weight model and an $\ell_1$ perturbation metric, pruning is predicted by the rate--distortion optimum itself.

\paragraph{Practical caveats.}
The clean perturbation bound above is stated for fully-connected networks with 1-Lipschitz activations.
Convolutional and transformer models require additional modeling choices, and in practice retraining or fine-tuning is often used after pruning.
Also, preserving accuracy is not the only possible criterion; compression can affect calibration, robustness, and fairness.
The information-theoretic view does not remove these engineering concerns, but it explains why sparse compressed reconstructions should be expected.

\subsubsection*{Connection to gradient compression}

The same Laplacian/extreme-value reasoning can also be applied to gradient compression in federated or distributed learning.
If gradient coordinates are approximately Laplacian, then communicating a small number of large residual coordinates is a rate--distortion motivated compression rule.
This resembles top-$k$ sparsification, but the threshold and coding interpretation are derived from the source distribution and the distortion objective.

% %%%%%%%%%%%%%%%%%%%%%%%%%%%%%%%%%%%%%%%%%%%%%%%%%%%%%%%%%%%%%%%%%%%%%%%%%%%%%%%
% % Lossy compression and rate--distortion
% %%%%%%%%%%%%%%%%%%%%%%%%%%%%%%%%%%%%%%%%%%%%%%%%%%%%%%%%%%%%%%%%%%%%%%%%%%%%%%%
% 
% \subsection{Setup: Rate--Distortion Tradeoff}
% \label{subsec:lossy-setup}
% 
% For real-valued sources, perfect reconstruction with a finite number of bits is impossible.
% A \emph{lossy} code trades rate (bits per symbol) for fidelity measured by a distortion function.
% 
% \paragraph{Single-letter distortion.}
% Let $X\sim P_X$ take values in $\cX$ (often $\cX\subseteq\fR^d$).
% Choose a distortion function $d:\cX\times\widehat{\cX}\to\fR_{\ge 0}$.
% For mean-squared error (MSE),
% \(
%     d(x,\hat x)=(x-\hat x)^2
% \)
% and the distortion level is $D=\E{(X-\widehat X)^2}$.
% 
% \paragraph{Block codes.}
% An $(n,R)$ lossy code consists of an encoder and decoder
% \[
%     f:\cX^n\to \{1,\dots,2^{nR}\},
%     \qquad
%     g:\{1,\dots,2^{nR}\}\to \widehat{\cX}^n,
% \]
% with reconstruction $\widehat X^n=g(f(X^n))$.
% We measure average per-letter distortion by
% \begin{equation}
%     d^{(n)}(x^n,\hat x^n)\coloneqq \frac{1}{n}\sum_{i=1}^n d(x_i,\hat x_i),
%     \qquad
%     \E{d^{(n)}(X^n,\widehat X^n)}.
% \end{equation}
% 
% \begin{definition}[Achievable $(R,D)$]
% A pair $(R,D)$ is \emph{achievable} if for every $\varepsilon>0$ and all sufficiently large $n$, there exists an $(n,R+\varepsilon)$ code such that
% \(\E{d^{(n)}(X^n,\widehat X^n)}\le D+\varepsilon\).
% \end{definition}
% 
% \begin{definition}[Rate--distortion function]
% The operational rate--distortion function is
% \begin{equation}
%     R(D)\coloneqq \inf\{R:\ (R,D)\text{ achievable}\}.
% \end{equation}
% \end{definition}
% 
% \subsection{Scalar Quantization: A Concrete Example}
% \label{subsec:scalar-quant}
% 
% As a warm-up, consider a scalar Gaussian source (e.g., heights).
% Assume $X\sim\cN(\mu,\sigma^2)$ with $\mu=170$ and $\sigma^2=10$.
% We want to store $X$ using $R$ bits per sample.
% 
% \paragraph{Rate $R=0$.}
% If we output a constant $\widehat X=\mu$, then
% \(D(0)=\Var{X}=\sigma^2\).
% 
% \paragraph{Rate $R=1$ (one threshold).}
% A 1-bit quantizer partitions at $\mu$:
% \(
%     f(x)=\I{x>\mu}.
% \)
% Given the index, the optimal reconstruction is the conditional mean,
% \(
%     \hat x_0=\E{X\mid X\le \mu},\ \hat x_1=\E{X\mid X>\mu}.
% \)
% For a Gaussian, this yields
% \begin{equation}
%     D(1)=\sigma^2\Bigl(1-\frac{2}{\pi}\Bigr)\approx 0.3634\,\sigma^2 \approx 3.63.
% \end{equation}
% 
% \paragraph{Rate $R=2$ (three thresholds).}
% With four cells (e.g., thresholds $165,170,175$) the distortion improves, but the choice of thresholds is not obviously optimal.
% 
% \subsection{Lloyd--Max Algorithm}
% \label{subsec:lloyd-max}
% 
% For a fixed number of reconstruction levels, the Lloyd--Max algorithm alternates between:
% \begin{enumerate}[leftmargin=2em]
%     \item \textbf{Centroid update:} with thresholds $\{\tau_i\}$ fixed,
%     \(\hat x_i \leftarrow \E{X\mid \tau_i\le X<\tau_{i+1}}\).
% 
%     \item \textbf{Boundary update:} with representatives $\{\hat x_i\}$ fixed,
%     \(\tau_{i+1}\leftarrow (\hat x_i+\hat x_{i+1})/2\) for squared-error distortion.
% \end{enumerate}
% Each step decreases distortion and the algorithm converges to a local optimum.
% 
% \begin{figure}[ht]
% \centering
%     \begin{subfigure}{0.40\textwidth}
%         \includegraphics[width=\linewidth]{scribe19_image2a.png}
%         \caption{Fix thresholds, update representatives}
%         \label{fig:lm-a}
%     \end{subfigure}
%     \hfill
%     \begin{subfigure}[b]{0.08\textwidth}
%         \includegraphics[width=\linewidth]{scribe19_arrow_both_directions.png}
%     \end{subfigure}
%     \hfill
%     \begin{subfigure}[b]{0.40\textwidth}
%         \includegraphics[width=\linewidth]{scribe19_image2b.png}
%         \caption{Fix representatives, update thresholds}
%         \label{fig:lm-b}
%     \end{subfigure}
% \caption{Alternating minimization in Lloyd--Max quantization.}
% \label{fig:lm}
% \end{figure}
% 
% \subsection{Vector Quantization}
% \label{subsec:vq}
% 
% Scalar quantization treats each dimension independently.
% Vector quantization (VQ) quantizes blocks $X^n$ jointly:
% \[
%   Q^{(n)} : \fR^{n} \longrightarrow \{1,\dots,2^{nR}\}.
% \]
% Even for independent components, VQ is strictly better at high rate because it can shape quantization regions to match the source geometry.
% 
% \begin{figure}[ht]
%   \centering
%   \includegraphics[width=.34\textwidth]{scribe19_image3.png}
%   \caption{Example 2-D source.}
%   \label{fig:vq-source}
% \end{figure}
% 
% \begin{figure}[ht]
%   \centering
%   \begin{subfigure}{0.41\textwidth}
%     \includegraphics[width=\linewidth]{scribe19_image4a.png}
%     \caption{Scalar quantization (rectangles).}
%     \label{fig:vq-scalar}
%   \end{subfigure}\hspace{8mm}
%   \begin{subfigure}{0.38\textwidth}
%     \includegraphics[width=\linewidth]{scribe19_image4b.png}
%     \caption{Vector quantization (adapted regions).}
%     \label{fig:vq-vector}
%   \end{subfigure}
%   \caption{Scalar vs. vector quantization for a 2-D uniform source.}
%   \label{fig:vq}
% \end{figure}
% 
% \begin{figure}[ht]
%   \centering
%   \includegraphics[width=0.30\textwidth]{scribe19_Figure3.jpg}
%   \caption{Rectangular cells incur extra MSE compared to more circular cells.}
%   \label{fig:vq-distortion}
% \end{figure}
% 
% \subsection{Information-Theoretic Rate--Distortion Function}
% \label{sec:rd-info}
% 
% The operational definition of $R(D)$ admits an exact single-letter characterization.
% 
% \begin{definition}[Information rate--distortion function]
% \label{def:rd-info}
% For a memoryless source $X\sim P_X$ and single-letter distortion $d$, define
% \begin{equation}
%     R^{(I)}(D)
%     \coloneqq
%     \inf_{P_{\widehat X\mid X}:\,\E{d(X,\widehat X)}\le D}
%     I(X;\widehat X).
% \end{equation}
% The infimum is over all conditional distributions (``test channels'') from $X$ to $\widehat X$.
% \end{definition}
% 
% \begin{theorem}[Shannon's rate--distortion theorem]
% \label{thm:shannon-rd}
% For every memoryless source and single-letter distortion measure,
% \begin{equation}
%     R(D)=R^{(I)}(D).
% \end{equation}
% \end{theorem}
% 
% \paragraph{Why mutual information appears.}
% A rate-$R$ code selects one out of $2^{nR}$ reconstructions.
% The random index cannot carry more than $nR$ bits about the source block.
% The theorem formalizes this by showing that the per-letter information that the reconstruction carries about the source is exactly $I(X;\widehat X)$.
% 
% \subsubsection*{Typicality lemma (intuition)}
% If we draw a random codeword $\widehat X^n$ i.i.d. from the marginal of a target joint distribution $P_{X,\widehat X}$, then the probability that $(X^n,\widehat X^n)$ is jointly typical (with respect to $P_{X,\widehat X}$) is roughly $2^{-nI(X;\widehat X)}$.
% Hence, to find \emph{one} jointly typical match among $2^{nR}$ codewords, we need $2^{nR}\cdot 2^{-nI}\gg 1$, i.e., $R>I(X;\widehat X)$.
% This is the heart of the random-coding achievability proof.
% 
% \subsection{Examples of Rate--Distortion Functions}
% \label{sec:rd-examples}
% 
% \subsubsection*{Binary source with Hamming distortion}
% 
% Let $X\sim\mathrm{Bern}(p)$ with $p\le \tfrac12$ and distortion
% \(d_H(x,\hat x)=\I{x\neq \hat x}\).
% Then
% \begin{equation}
% R(D)=
% \begin{cases}
% \Hb{p}-\Hb{D}, & 0\le D\le p,\\
% 0, & D\ge p.
% \end{cases}
% \end{equation}
% 
% \paragraph{Sketch.}
% Let $Z\coloneqq X\oplus \widehat X$ be the error bit.
% Then $H(X\mid \widehat X)=H(Z\mid \widehat X)\le H(Z)\le \Hb{D}$.
% So $I(X;\widehat X)=H(X)-H(X\mid \widehat X)\ge \Hb{p}-\Hb{D}$.
% Achievability uses a binary symmetric test channel with crossover $D$.
% 
% \subsubsection*{Gaussian source with quadratic distortion}
% 
% Let $X\sim\cN(0,\sigma^2)$ and $d(x,\hat x)=(x-\hat x)^2$.
% Then
% \begin{equation}
% R(D)=
% \begin{cases}
% \tfrac12\log\Bigl(\frac{\sigma^2}{D}\Bigr), & 0<D\le \sigma^2,\\
% 0, & D\ge \sigma^2.
% \end{cases}
% \end{equation}
% 
% \paragraph{Sketch.}
% For any reconstruction $\widehat X$ with $\E{(X-\widehat X)^2}\le D$,
% \[
% I(X;\widehat X)=h(X)-h(X\mid \widehat X)=h(X)-h(X-\widehat X\mid \widehat X)
% \ge h(X)-h(X-\widehat X).
% \]
% By Theorem~\ref{thm:gauss-maxent}, $h(X-\widehat X)\le \tfrac12\log(2\pi e D)$, so
% \(I(X;\widehat X)\ge \tfrac12\log(\sigma^2/D)\).
% Achievability uses a Gaussian test channel
% \(X=\widehat X+Z\) with $Z\sim\cN(0,D)$ independent.
% 
% \paragraph{Geometric picture.}
% In $\fR^n$, a typical Gaussian source block lies near a sphere of radius $\sqrt{n\sigma^2}$.
% A distortion constraint $D$ means each codeword covers a ball of radius $\sqrt{nD}$.
% Comparing volumes yields the same bound $R\ge \tfrac12\log(\sigma^2/D)$.
% 
% \begin{figure}[h]
% \centering
% \includegraphics[width=0.58\textwidth]{scribe27_sphere.jpg}
% \caption{High-dimensional picture for Gaussian rate--distortion: cover a large sphere (source) by many small balls (distortion).}
% \label{fig:gauss-geometry}
% \end{figure}
% 
% \subsubsection*{Laplacian source with absolute distortion}
% 
% Let $X\sim\mathrm{Lap}(1/\lambda)$ with pdf
% \(f_X(x)=\tfrac{\lambda}{2}e^{-\lambda|x|}\)
% and distortion $d_1(x,\hat x)=|x-\hat x|$.
% Then
% \begin{equation}
% R(D)=
% \begin{cases}
% -\log(\lambda D), & 0<D\le 1/\lambda,\\
% 0, & D\ge 1/\lambda.
% \end{cases}
% \end{equation}
% 
% \paragraph{One useful consequence: sparsity in the optimal reproduction.}
% A test channel that achieves the bound has a reproduction distribution of the form
% \begin{equation}
%     f_{\widehat X}(\hat x)
%     = \lambda^2D^2\,\delta(\hat x)
%     + (1-\lambda^2D^2)\,\frac{\lambda}{2}e^{-\lambda|\hat x|},
% \end{equation}
% which has an atom at zero.
% Thus, \emph{rate--distortion optimal reconstructions are intrinsically sparse} for Laplacian-like sources under $\ell_1$ distortion.
% This observation will reappear in neural network pruning.
% 
% \subsection{Extreme-Value Theory (Tool for Sparse Refinement)}
% \label{sec:evt}
% 
% In some iterative compression schemes we repeatedly identify unusually large coordinates of a residual vector.
% Extreme-value theory tells us where these maxima typically lie.
% 
% \begin{example}[Max of exponentials]
% If $Y_1,\dots,Y_n\overset{\text{i.i.d.}}{\sim}\mathrm{Exp}(\lambda)$ and $M_n=\max_i Y_i$, then
% \begin{equation}
%     M_n = \frac{\log n}{\lambda} + O_p(1).
% \end{equation}
% The fluctuations converge to a Gumbel distribution after centering by $\frac{1}{\lambda}\log n$.
% \end{example}
% 
% \paragraph{Rule of thumb.}
% For i.i.d. $X_i\sim\cN(0,\sigma^2)$,
% \(\max_i X_i = \sigma\sqrt{2\log n}+O_p(1)\).
% 
% \subsection{Application: Compressing Neural Networks}
% \label{sec:nn-compress}
% 
% Modern neural network (NN) compression often starts by choosing a distortion metric on weights that upper-bounds the induced change in the network output.
% Rate--distortion theory then provides a principled limit for how well such weights can be compressed.
% 
% \subsubsection*{A distortion metric that controls output perturbation}
% 
% Consider a $d$-layer fully-connected network
% \(f(x;w)=W^{(d)}\sigma(\cdots\sigma(W^{(1)}x)\cdots)\)
% with a 1-Lipschitz nonlinearity $\sigma$ (e.g., ReLU).
% A typical bound is:
% 
% \begin{theorem}[Weight distortion implies output distortion]
% \label{thm:nn-l1-bound}
% Assume $\|\widehat W^{(\ell)}\|_1\le \|W^{(\ell)}\|_1$ for all layers.
% Then
% \begin{equation}
% \sup_{\|x\|_1\le 1}\,\|f(x;w)-f(x;\widehat w)\|_1
% \le
% \Bigl(\prod_{k=1}^{d}\|W^{(k)}\|_1\Bigr)
% \Bigl(\sum_{\ell=1}^{d}\frac{\|W^{(\ell)}-\widehat W^{(\ell)}\|_1}{\|W^{(\ell)}\|_1}\Bigr).
% \end{equation}
% \end{theorem}
% 
% This motivates a normalized $\ell_1$ distortion on weights.
% Empirically, many pretrained weight distributions are close to Laplacian after normalization.
% Then the Laplacian $\ell_1$ rate--distortion calculation above predicts that an optimal compressor should generate \emph{sparse} reconstructions, providing an information-theoretic justification for pruning.
% 
% \subsubsection*{Successive refinement viewpoint}
% 
% A direct random codebook achieving the rate--distortion function is computationally infeasible for models with millions of weights.
% One practical alternative is to iteratively encode the residual error.
% If the residual after an optimal step has the same distributional form as the original source (\emph{successive refinability}), then repeated refinement can approach the information-theoretic limit with manageable complexity.
% 